\section{Experimental Setup}

All benchmark programs were implemented in C++ and compiled using \texttt{g++} with optimization level \texttt{-O3}. Runtime measurements were obtained using the C++ \texttt{std::chrono} high-resolution clock. Each benchmark was executed repeatedly in order to reduce noise caused by background system activity and operating system scheduling \cite{georges2007}.

Input datasets for the sorting experiments were generated using a Python script and written to disk as plain text files. Three dataset sizes were created to evaluate algorithmic scaling, along with three ordering configurations: random order, pre-sorted order, and reverse-sorted order.

The benchmark program loaded each dataset into memory and measured the execution time of the sorting algorithm. Each dataset was tested multiple times, and the resulting runtime measurements were recorded in CSV files for later analysis.

A second benchmark was implemented to evaluate memory locality. This experiment performed traversal of large matrices using two access patterns: row-major traversal and column-major traversal. Both approaches perform the same number of arithmetic operations but differ in how memory is accessed. Because modern CPUs rely heavily on cache hierarchies, sequential memory access patterns are expected to exhibit better performance due to improved spatial locality \cite{lam1991, mckee2004}.
